% Noether P09 Korean producer draft — UNCHECKED
% T01-U07; ED0002 lines6431--6446; source 1,004 LF B / EEC6E00E924F0AFC86C5A9B9AF104297B3029AA1F2EBE9B1DA56C05DF953B191
% Pointer v009 / ED0002: B06BE3530D9CF2E82B56FDBA7FE41D5D044DF2425DFA2A059D4939EAA2F7A6C2 / C9A125167ACB33D914EE4374B65AE7CDF0052F568371B8B77B720EA178ABF0E3
% Translation only; no review, build, render, or approval.

I--IV와 더불어 V도 만족하는 특수한 영역을 \emph{대수적으로 정수인 초월수}의 영역 $\mathfrak H$라 부르자. 가장 일반적인 영역 $\mathfrak H$에 대해서는 다음의 단순한 결과가 얻어진다(§ 4).
\begin{enumerate}
\item[a)] \emph{같은 기저 $H$를 이용하여 구성할 수 있는 모든 영역 $\mathfrak H$의 교집합은 그 기저의 모든 정수적 대수함수로 주어진다. 따라서 그 교집합은 체르멜로 영역 $\mathfrak G_\eta$이며, 이 영역 자체도 영역 $\mathfrak H$이다.} (이로써 동시에 체르멜로 영역의 추상적 특징짓기도 얻어진다.)
\item[b)] \emph{모든 영역 $\mathfrak H$는 임의의 계 $\mathfrak S(\eta)$를 $\mathfrak G_\eta$에 대수적으로 정수인 방식으로 첨가하여 얻어진다.\srcfn{**)}{이 표현의 설명은 § 4에 있다.} 여기서 $\mathfrak S(\eta)$가 만족해야 하는 유일한 조건은 그 첨가로 어떠한 분수형 대수수도 그 영역 안으로 들어오지 않는다는 것이다.}
\end{enumerate}
